\subsection{Zahlenmengen}
\label{subsec:fundamentals_number_sets}

Ausgehend von den natürlichen Zahlen werden die Zahlenmengen schrittweise
erweitert. Der Grund ist stets derselbe: Eine Rechenoperation soll ausführbar
sein, ihr Ergebnis liegt aber nicht mehr in der bisherigen Zahlenmenge.

\begin{DefinitionBox}{Abgeschlossenheit}{fundamentals_closure}
  Eine Zahlenmenge $M$ heißt bezüglich einer Rechenoperation $\circ$
  \emph{abgeschlossen} \EnglishTerm{closed}, wenn für alle $a,b\in M$ auch
  \[
    a\circ b\in M
  \]
  gilt, sofern die Rechenoperation definiert ist.
\end{DefinitionBox}

\subsubsection{Natürliche Zahlen}
\label{subsubsec:fundamentals_natural_numbers}

Die \emph{natürlichen Zahlen} dienen dem Zählen:
\[
  \N=\{1,2,3,4,\ldots\},
  \qquad
  \N_0=\N\cup\{0\}=\{0,1,2,3,\ldots\}.
\]
Die Addition und Multiplikation natürlicher Zahlen führen wieder zu einer
natürlichen Zahl. Die natürlichen Zahlen sind bezüglich $+$ und $\cdot$
abgeschlossen:
\[
  3+5=8\in\N,
  \qquad
  3\cdot5=15\in\N.
\]

Bezüglich der Subtraktion sind die natürlichen Zahlen dagegen nicht
abgeschlossen:
\[
  3-5=-2\notin\N.
\]

\subsubsection{Ganze Zahlen}
\label{subsubsec:fundamentals_integer_numbers}

Damit auch solche Differenzen erfasst werden, ergänzt man die negativen Zahlen
und erhält die \emph{ganzen Zahlen}:
\[
  \Z=\{\ldots,-3,-2,-1,0,1,2,3,\ldots\}.
\]
In $\Z$ sind Addition, Subtraktion und Multiplikation abgeschlossen.

\begin{center}
  \begingroup
    \def\NumberSetStage{2}
    \begin{tikzpicture}[x=1cm,y=1cm]
  \ifnum\NumberSetStage>4\relax
    \fill[THMBlue!34] (0,0) ellipse (7.5cm and 2.6cm);
    \draw[THMBlue!85!black,line width=1pt] (0,0) ellipse (7.5cm and 2.6cm);
  \fi

  \ifnum\NumberSetStage>3\relax
    \fill[THMBlue!25] (0,0) ellipse (6.1cm and 2.1cm);
    \draw[THMBlue!72!black,line width=1pt] (0,0) ellipse (6.1cm and 2.1cm);
  \fi

  \ifnum\NumberSetStage>2\relax
    \fill[THMBlue!17] (0,0) ellipse (4.7cm and 1.6cm);
    \draw[THMBlue!60!black,line width=1pt] (0,0) ellipse (4.7cm and 1.6cm);
  \fi

  \fill[THMBlue!10] (0,0) ellipse (3.3cm and 1.1cm);
  \draw[THMBlue!48!black,line width=1pt] (0,0) ellipse (3.3cm and 1.1cm);

  \fill[THMBlue!4] (0,0) ellipse (1.8cm and 0.55cm);
  \draw[THMBlue!35!black,line width=1pt] (0,0) ellipse (1.8cm and 0.55cm);

  \node[font=\Large\bfseries] at (0,0) {$\N$};
  \node[font=\Large\bfseries] at (2.55,0) {$\Z$};
  \ifnum\NumberSetStage>2\relax
    \node[font=\Large\bfseries] at (4.0,0) {$\Q$};
  \fi
  \ifnum\NumberSetStage>3\relax
    \node[font=\Large\bfseries] at (5.4,0) {$\R$};
  \fi
  \ifnum\NumberSetStage>4\relax
    \node[font=\Large\bfseries] at (6.8,0) {$\C$};
  \fi
\end{tikzpicture}

  \endgroup
\end{center}

Bezüglich der Division sind die ganzen Zahlen nicht abgeschlossen:
\[
  3:5=\frac35=0{,}6\notin\Z.
\]

\subsubsection{Rationale Zahlen}
\label{subsubsec:fundamentals_rational_numbers}

Damit auch Quotienten ganzer Zahlen erfasst werden, erweitert man $\Z$ zu den
\emph{rationalen Zahlen}:
\[
  \Q=\left\{\frac pq\,\middle|\,p\in\Z,\ q\in\N\right\}.
\]
Jede rationale Zahl lässt sich als Bruch ganzer Zahlen darstellen. Endliche
und periodische Dezimalzahlen sind daher rational, zum Beispiel
\[
  -\frac12=-0{,}5\in\Q,
  \qquad
  \frac19=0{,}\overline{1}\in\Q.
\]
Die rationalen Zahlen sind bezüglich $+$, $-$, $\cdot$ und der Division durch
eine von null verschiedene Zahl abgeschlossen.

\begin{center}
  \begingroup
    \def\NumberSetStage{3}
    \begin{tikzpicture}[x=1cm,y=1cm]
  \ifnum\NumberSetStage>4\relax
    \fill[THMBlue!34] (0,0) ellipse (7.5cm and 2.6cm);
    \draw[THMBlue!85!black,line width=1pt] (0,0) ellipse (7.5cm and 2.6cm);
  \fi

  \ifnum\NumberSetStage>3\relax
    \fill[THMBlue!25] (0,0) ellipse (6.1cm and 2.1cm);
    \draw[THMBlue!72!black,line width=1pt] (0,0) ellipse (6.1cm and 2.1cm);
  \fi

  \ifnum\NumberSetStage>2\relax
    \fill[THMBlue!17] (0,0) ellipse (4.7cm and 1.6cm);
    \draw[THMBlue!60!black,line width=1pt] (0,0) ellipse (4.7cm and 1.6cm);
  \fi

  \fill[THMBlue!10] (0,0) ellipse (3.3cm and 1.1cm);
  \draw[THMBlue!48!black,line width=1pt] (0,0) ellipse (3.3cm and 1.1cm);

  \fill[THMBlue!4] (0,0) ellipse (1.8cm and 0.55cm);
  \draw[THMBlue!35!black,line width=1pt] (0,0) ellipse (1.8cm and 0.55cm);

  \node[font=\Large\bfseries] at (0,0) {$\N$};
  \node[font=\Large\bfseries] at (2.55,0) {$\Z$};
  \ifnum\NumberSetStage>2\relax
    \node[font=\Large\bfseries] at (4.0,0) {$\Q$};
  \fi
  \ifnum\NumberSetStage>3\relax
    \node[font=\Large\bfseries] at (5.4,0) {$\R$};
  \fi
  \ifnum\NumberSetStage>4\relax
    \node[font=\Large\bfseries] at (6.8,0) {$\C$};
  \fi
\end{tikzpicture}

  \endgroup
\end{center}

Bezüglich des Wurzelziehens sind die rationalen Zahlen nicht abgeschlossen. So
gibt es keine rationale Zahl, deren Quadrat gleich $2$ ist:
\[
  x^2=2
  \quad\Longrightarrow\quad
  x=\sqrt2\notin\Q.
\]

\subsubsection{Reelle Zahlen}
\label{subsubsec:fundamentals_real_numbers}

Zahlen, die sich nicht als Quotient ganzer Zahlen darstellen lassen, heißen
\emph{irrationale Zahlen}. Beispiele sind
\[
  \sqrt2=1{,}4142\ldots,
  \qquad
  \pi=3{,}14159\ldots,
  \qquad
  \mathrm e=2{,}71828\ldots.
\]
Rationale und irrationale Zahlen bilden zusammen die \emph{reellen Zahlen}.

\begin{center}
  \begingroup
    \def\NumberSetStage{4}
    \begin{tikzpicture}[x=1cm,y=1cm]
  \ifnum\NumberSetStage>4\relax
    \fill[THMBlue!34] (0,0) ellipse (7.5cm and 2.6cm);
    \draw[THMBlue!85!black,line width=1pt] (0,0) ellipse (7.5cm and 2.6cm);
  \fi

  \ifnum\NumberSetStage>3\relax
    \fill[THMBlue!25] (0,0) ellipse (6.1cm and 2.1cm);
    \draw[THMBlue!72!black,line width=1pt] (0,0) ellipse (6.1cm and 2.1cm);
  \fi

  \ifnum\NumberSetStage>2\relax
    \fill[THMBlue!17] (0,0) ellipse (4.7cm and 1.6cm);
    \draw[THMBlue!60!black,line width=1pt] (0,0) ellipse (4.7cm and 1.6cm);
  \fi

  \fill[THMBlue!10] (0,0) ellipse (3.3cm and 1.1cm);
  \draw[THMBlue!48!black,line width=1pt] (0,0) ellipse (3.3cm and 1.1cm);

  \fill[THMBlue!4] (0,0) ellipse (1.8cm and 0.55cm);
  \draw[THMBlue!35!black,line width=1pt] (0,0) ellipse (1.8cm and 0.55cm);

  \node[font=\Large\bfseries] at (0,0) {$\N$};
  \node[font=\Large\bfseries] at (2.55,0) {$\Z$};
  \ifnum\NumberSetStage>2\relax
    \node[font=\Large\bfseries] at (4.0,0) {$\Q$};
  \fi
  \ifnum\NumberSetStage>3\relax
    \node[font=\Large\bfseries] at (5.4,0) {$\R$};
  \fi
  \ifnum\NumberSetStage>4\relax
    \node[font=\Large\bfseries] at (6.8,0) {$\C$};
  \fi
\end{tikzpicture}

  \endgroup
\end{center}

Bezüglich des Wurzelziehens sind auch die reellen Zahlen nicht abgeschlossen:
Die Gleichung
\[
  x^2=-1
\]
besitzt in $\R$ keine Lösung, denn das Quadrat jeder reellen Zahl ist
nichtnegativ.

\subsubsection{Komplexe Zahlen}
\label{subsubsec:fundamentals_complex_numbers}

Man führt die \emph{imaginäre Einheit} $\mathrm i$ mit
\[
  \mathrm i^2=-1
\]
ein. Die dadurch entstehenden \emph{komplexen Zahlen} haben die Form
\[
  \C=\{a+b\mathrm i\mid a,b\in\R\}.
\]
Jede reelle Zahl ist wegen $a=a+0\mathrm i$ zugleich eine komplexe Zahl. In
$\C$ sind Addition, Subtraktion, Multiplikation und die Division durch eine von
null verschiedene Zahl abgeschlossen; außerdem ist nun $x^2=-1$ lösbar:
\[
  x=\pm\mathrm i\in\C.
\]

\begin{center}
  \begingroup
    \def\NumberSetStage{5}
    \begin{tikzpicture}[x=1cm,y=1cm]
  \ifnum\NumberSetStage>4\relax
    \fill[THMBlue!34] (0,0) ellipse (7.5cm and 2.6cm);
    \draw[THMBlue!85!black,line width=1pt] (0,0) ellipse (7.5cm and 2.6cm);
  \fi

  \ifnum\NumberSetStage>3\relax
    \fill[THMBlue!25] (0,0) ellipse (6.1cm and 2.1cm);
    \draw[THMBlue!72!black,line width=1pt] (0,0) ellipse (6.1cm and 2.1cm);
  \fi

  \ifnum\NumberSetStage>2\relax
    \fill[THMBlue!17] (0,0) ellipse (4.7cm and 1.6cm);
    \draw[THMBlue!60!black,line width=1pt] (0,0) ellipse (4.7cm and 1.6cm);
  \fi

  \fill[THMBlue!10] (0,0) ellipse (3.3cm and 1.1cm);
  \draw[THMBlue!48!black,line width=1pt] (0,0) ellipse (3.3cm and 1.1cm);

  \fill[THMBlue!4] (0,0) ellipse (1.8cm and 0.55cm);
  \draw[THMBlue!35!black,line width=1pt] (0,0) ellipse (1.8cm and 0.55cm);

  \node[font=\Large\bfseries] at (0,0) {$\N$};
  \node[font=\Large\bfseries] at (2.55,0) {$\Z$};
  \ifnum\NumberSetStage>2\relax
    \node[font=\Large\bfseries] at (4.0,0) {$\Q$};
  \fi
  \ifnum\NumberSetStage>3\relax
    \node[font=\Large\bfseries] at (5.4,0) {$\R$};
  \fi
  \ifnum\NumberSetStage>4\relax
    \node[font=\Large\bfseries] at (6.8,0) {$\C$};
  \fi
\end{tikzpicture}

  \endgroup
\end{center}
